\documentclass[10pt]{mypackage}

\usepackage[uprightgreeks,light]{kpfonts}
\renewcommand{\coloneq}{\coloneqq}
%\renewcommand{\mathbb}{\mathds}
\usepackage{homework}
%\usepackage{notes}

%\usepackage[ backend=bibtex, style = alphabetic, sorting=ynt ]{biblatex}
%\addbibresource{  }

\usepackage{parskip}

\fancyhf{}
\fancyhead[R]{Avinash Iyer}
\fancyhead[L]{Algebraic Topology: Homework 24 Revisions}
\fancyfoot[C]{\thepage}

\setcounter{secnumdepth}{0}

\begin{document}
\RaggedRight
\section{Revised Problems}%
\begin{problem}[Problem 1]\hfill
  \begin{enumerate}[(a)]
    \item Prove that $ \widetilde{H}_q\left( X,A \right) \cong H_q\left( X,A \right)$ for all $q\geq 0$.
    \item Let $x_0\in X$. Then, $ \widetilde{H}_q\left( X \right)\cong H_q\left( X,\set{x_0} \right) $.
  \end{enumerate}
\end{problem}
\begin{solution}\hfill
  \begin{enumerate}[(a)]
    \item For any $q\geq 1$, we have that $ \widetilde{C}_q\left( X \right) = C_q(X) $ and $ \widetilde{C}_q(A) = C_q(A) $ with identical boundary maps at each step excluding level $0$. The relative cycles in $ \widetilde{C}_q\left( X \right) $ are then equal to the relative cycles in $ C_q(X) $, and have identical relative boundaries, meaning that for any $q\geq 1$, we have $\widetilde{H}_q\left( X,A \right) = H_q\left( X,A \right)$ since their definitions and elements are exactly the same at these points.

      At level $0$, the fact that both the augmentation homomorphism and the boundary map $\partial_1$ descend to $A$ gives that
      \begin{align*}
        \widetilde{H}_0\left( X \right) &= \frac{\ker\left( \ve \right)}{\img\left( \partial_1 \right)}\\
        \widetilde{H}_0\left( A \right) &= \frac{\ker\left( \ve \right)\cap C_0\left( A \right)}{\img\left( \partial_1 \right)\cap C_0\left( A \right)},
      \end{align*}
      while the fact that relative boundaries are given by $ \img\left( \partial_1 \right) + \widetilde{C}_0\left( A \right) $ gives
      \begin{align*}
        \widetilde{H}_0\left( X,A \right) &= \frac{\ker\left( \ve \right)}{ \img\left( \partial_1 \right) + \left( \ker\left( \ve \right)\cap C_0\left( A \right) \right) }\\
                                          &= \frac{\ker\left( \ve \right)}{\ker\left( \ve \right)\cap \left( \img\left( \partial_1 \right) + C_0\left( A \right) \right)}.
      \end{align*}
      From the second isomorphism theorem, we have
      \begin{align*}
        \frac{\ker\left( \ve \right)}{\ker\left( \ve \right)\cap \left( \img\left( \partial_1 \right) + C_0\left( A \right) \right)} &\cong \frac{\ker\left( \ve \right) + \img\left( \partial_1 \right) + C_0\left( A \right)}{\img\left( \partial_1 \right) + C_0\left( A \right)}.
      \end{align*}
      Unfortunately, I think I might have lost some definitions somewhere here that prevent me from achieving my objective, which is to show that this eventually equals $\frac{C_0\left( X \right)}{\img\left( \partial_1 \right) + C_0\left( A \right)} = H_0\left( X,A \right)$.
    \item We have that $ H_q\left( X,\set{x_0} \right)\cong \widetilde{H}_q\left( X,\set{x_0} \right) $ from part (a). In the long exact sequence of the pair $\left( \set{x_0},X \right)$, we thus have
      \begin{align*}
        \widetilde{H}_q\left( \set{x_0} \right) \rightarrow \widetilde{H}_q\left( X \right) \rightarrow \widetilde{H}_q\left( X,\set{x_0} \right) \rightarrow \widetilde{H}_{q-1}\left( \set{x_0} \right).
      \end{align*}
      Since $ \widetilde{H}_q\left( \set{x_0} \right) = 0 $ for all $q$, it follow that $ \widetilde{H}_q\left( X \right)\cong \widetilde{H}_q\left( X,\set{x_0} \right) $, whence $ \widetilde{H}_q\left( X \right)\cong H_q\left( X,\set{x_0} \right) $.
  \end{enumerate}
  
\end{solution}
\end{document}
